\chapter{RSA Public-Key Cryptography}

RSA (Rivest, Shamir, Adleman, 1977) is the payoff of the number theory so far: it uses the
Euclidean algorithm, modular inverses, Fermat's Little Theorem and the splitting modulus
theorem.

\section{Public-key cryptography}

In \emph{symmetric} (private-key) cryptography, sender and receiver share a secret key used
both to encrypt and decrypt; the problem is agreeing on that key securely. In
\emph{public-key} cryptography, each receiver has two keys:
\begin{itemize}
  \item a \emph{public key}, published to everyone and used to \emph{encrypt} messages to her;
  \item a \emph{private key}, kept secret and used to \emph{decrypt}.
\end{itemize}
Security rests on the belief that computing the private key from the public key is
infeasible. For RSA, this is (essentially) the difficulty of factoring a large $n = pq$.

\section{The RSA scheme}

\begin{important}[RSA key generation]
\begin{enumerate}
  \item Choose two distinct large primes $p$ and $q$ (hundreds of digits in practice).
  \item Compute $n = pq$ and $\varphi = (p-1)(q-1)$.
  \item Choose $e$ with $1 < e < \varphi$ and $\gcd(e,\varphi) = 1$.
  \item Compute $d$ with $1 < d < \varphi$ and $ed \equiv 1 \Mod{\varphi}$ (EEA).
  \item Public key: $(e, n)$. Private key: $(d, n)$. Keep $p$, $q$, $\varphi$ secret.
\end{enumerate}
\end{important}

\begin{important}[Encryption and decryption]
A message is an integer $M$ with $0 \le M < n$.
\begin{itemize}
  \item Encrypt: $C \equiv M^e \Mod n$, with $0 \le C < n$.
  \item Decrypt: $R \equiv C^d \Mod n$, with $0 \le R < n$. Then $R = M$.
\end{itemize}
\end{important}

Why it is (believed) secure: an eavesdropper knows $n$, $e$ and $C$. To find $d$ she needs
$\varphi$; but knowing $n$ and $\varphi = n - (p+q) + 1$ is equivalent to knowing $p + q$ and
$pq$, hence $p$ and $q$ (roots of $t^2 - (p+q)t + pq$). So recovering $d$ this way is as hard as
factoring $n$.

All the steps for the legitimate users are fast: primality testing, the EEA, and modular
exponentiation by repeated squaring each take a number of steps polynomial in the number of
digits.

\section{A small worked example}

\begin{example}[Key generation]
Take $p = 11$, $q = 13$. Then $n = 143$ and $\varphi = 10 \cdot 12 = 120$. Choose $e = 7$;
$\gcd(7, 120) = 1$ since $120 = 2^3\cdot 3\cdot 5$. To find $d$, run the EEA on $120$ and $7$:
\begin{center}
\begin{tabular}{rrrr}
\toprule
$x$ & $y$ & $r$ & $q$\\
\midrule
$1$ & $0$ & $120$ & \\
$0$ & $1$ & $7$ & \\
$1$ & $-17$ & $1$ & $17$\\
\bottomrule
\end{tabular}
\end{center}
So $120(1) + 7(-17) = 1$, i.e.\ $7\cdot(-17) \equiv 1 \Mod{120}$, and
$d \equiv -17 \equiv 103 \Mod{120}$. Check: $7 \cdot 103 = 721 = 6\cdot 120 + 1$.

Public key $(e,n) = (7, 143)$; private key $(d,n) = (103, 143)$.
\end{example}

\begin{example}[Encryption]
Encrypt $M = 9$: compute $9^7 \bmod 143$ by repeated squaring.
\begin{align*}
9^2 &= 81,\\
9^4 &= 81^2 = 6561 = 45\cdot 143 + 126 \equiv 126 \equiv -17,\\
9^7 &= 9^4\cdot 9^2\cdot 9 \equiv (-17)(81)(9) = (-17)(729).
\end{align*}
Since $729 = 5\cdot 143 + 14$, this is $(-17)(14) = -238 \equiv -238 + 2\cdot 143 = 48$.
So the ciphertext is $C = 48$.
\end{example}

\begin{example}[Decryption]
Compute $48^{103} \bmod 143$. The receiver knows $p,q$, so she may use SMT and FLT.

\emph{Mod $11$:} $48 \equiv 4$. By FLT exponents reduce mod $10$: $103 \equiv 3$, so
$48^{103} \equiv 4^3 = 64 \equiv 9 \Mod{11}$.

\emph{Mod $13$:} $48 \equiv 9$. Exponents reduce mod $12$: $103 = 8\cdot 12 + 7$, so we need
$9^7$. $9^2 = 81 \equiv 3$, $9^4 \equiv 9$, $9^7 = 9^4\cdot 9^2\cdot 9 \equiv 9 \cdot 3\cdot 9 = 243 =
18\cdot 13 + 9\equiv 9 \Mod{13}$.

So $R \equiv 9 \Mod{11}$ and $R \equiv 9 \Mod{13}$; by CRT/SMT, $R \equiv 9 \Mod{143}$. We recover
$M = 9$. (Without knowing $p,q$, the same result follows from repeated squaring mod $143$
directly, just with more work.)
\end{example}

\section{Proof of correctness}

\begin{theorem}[RSA works]
With $n, e, d$ as above and $0 \le M < n$, if $C \equiv M^e \Mod n$ and $R \equiv C^d\Mod n$
with $0 \le R < n$, then $R = M$.
\end{theorem}

\begin{proof}
Since $ed \equiv 1 \Mod{\varphi}$, write $ed = 1 + k(p-1)(q-1)$ for some integer $k\ge 0$. Then
\[ R \equiv C^d \equiv M^{ed} \Mod n. \]
We show $M^{ed} \equiv M \Mod p$.
\begin{itemize}
  \item If $p \nmid M$: by FLT $M^{p-1} \equiv 1 \Mod p$, so
        $M^{ed} = M\cdot\bigl(M^{p-1}\bigr)^{k(q-1)} \equiv M\cdot 1 = M \Mod p$.
  \item If $p \mid M$: both $M^{ed}$ and $M$ are $\equiv 0 \Mod p$.
\end{itemize}
By the same argument, $M^{ed} \equiv M \Mod q$. Since $p \neq q$ are primes,
$\gcd(p,q) = 1$, and SMT gives $M^{ed} \equiv M \Mod{pq}$. So $R \equiv M \Mod n$, and since both
$R$ and $M$ lie in $\{0, 1, \dots, n-1\}$, CISR gives $R = M$.
\end{proof}

\begin{important}[Why $p \ne q$ and $\gcd(e,\varphi) = 1$]
$\gcd(e,\varphi) = 1$ is exactly what makes the inverse $d$ exist (INV $\Zm m$). Distinctness
of $p$ and $q$ is used when we apply SMT; with $p = q$, $n = p^2$ is trivially factored
anyway.
\end{important}

\section{Digital signatures}

RSA also gives signatures. To sign $M$, Alice computes $S \equiv M^{d} \Mod{n}$ with her
\emph{private} key. Anyone can verify by checking $S^{e} \equiv M \Mod n$ with her public key;
only someone knowing $d$ could have produced $S$. Correctness is the same theorem with the roles
of $e$ and $d$ swapped (since $ed = de$).

\begin{example}
With Alice's keys above, signing $M = 2$: $S \equiv 2^{103} \Mod{143}$. Mod $11$: $2^{3} = 8$.
Mod $13$: $2^{7} = 128 = 9\cdot 13 + 11 \equiv 11$. Solve $S\equiv 8 \Mod{11}$,
$S \equiv 11 \Mod{13}$: $S = 11 + 13k$, $11 + 2k \equiv 8 \Mod{11}$, so $2k \equiv 8$ and
$k \equiv 4$, giving $S = 63$. Verify: $63^7 \bmod 143$. Mod $11$: $63 \equiv 8$, and
$8^7 \equiv 8^{-3} \equiv 2$ (as $8^3 = 512 = 46\cdot 11 + 6$ and $6\cdot 2 = 12 \equiv 1$).
Mod $13$: $63 \equiv 11 \equiv -2$, $(-2)^7 = -128 \equiv -128 + 130 = 2$. So $63^7 \equiv 2$,
as required.
\end{example}

\section{Common pitfalls}
\begin{itemize}
  \item Inverting $e$ modulo $n$ instead of modulo $\varphi = (p-1)(q-1)$.
  \item Choosing $e$ with $\gcd(e,\varphi) \ne 1$ (e.g.\ any even $e$, since $\varphi$ is even).
  \item Forgetting the case $p \mid M$ in the correctness proof.
  \item Arithmetic errors in repeated squaring; reduce after every multiplication and use
        negative representatives to keep numbers small.
\end{itemize}

\section{Summary}
\begin{itemize}
  \item Keys: $n = pq$, $\varphi = (p-1)(q-1)$, $\gcd(e,\varphi) = 1$, $ed \equiv 1 \Mod \varphi$;
        public $(e,n)$, private $(d,n)$.
  \item $C \equiv M^e$, $M \equiv C^d \Mod n$.
  \item Correctness: FLT modulo $p$ and $q$ separately, then SMT.
  \item Security relies on the difficulty of factoring $n$; legitimate operations are fast.
\end{itemize}
